\chapter{Stato dell’arte}
\label{ch:sota}

La guida cooperativa in \emph{platooning} mira a migliorare sicurezza, capacità stradale ed efficienza energetica tramite il coordinamento longitudinale di più veicoli in colonna. Per ottenere distanze inter-veicolari ridotte in modo sicuro, i sistemi di controllo devono reagire rapidamente alle variazioni di velocità e alle perturbazioni, e spesso beneficiano dello scambio di informazioni tra veicoli oltre ai soli sensori locali. Questo rende la qualità delle comunicazioni un elemento strutturale del problema: ritardi, perdite e discontinuità del link si riflettono direttamente su stabilità e vincoli di sicurezza, e impongono un co-design tra rete e controllo \cite{dressler2019cooperative,giordano2019joint}.

Negli ultimi anni è emerso con chiarezza che nessuna singola tecnologia radio è in grado di garantire, in ogni condizione operativa, l’affidabilità richiesta dalle applicazioni safety-critical del platooning. Fenomeni come congestione, shadowing, \emph{blind spot} e variazioni repentine del canale possono compromettere comunicazioni intraplatoon sia su tecnologie dedicate (es.\ IEEE~802.11p) sia su tecnologie cellulari, con impatti misurabili sulla capacità del sistema di mantenere prestazioni e sicurezza \cite{segata2015communication,bazzi2020wireless}. Per questo motivo stanno prendendo forma approcci multi-tecnologia (ad es.\ combinazioni radio/VLC e radio/cellulare) che sfruttano \emph{diversity} e complementarità tra canali, ma che introducono nuove complessità: selezione del link, ridondanza, gestione di \emph{fallback} e \emph{recovery}, e transizioni sicure tra controllori in funzione della connettività disponibile \cite{ishihara2015improving,segata2016platooning,segata2023multi-technology}.

Questo capitolo fornisce il contesto scientifico e tecnico necessario per inquadrare tali problematiche e motivare le scelte progettuali e sperimentali adottate nel seguito. In particolare, vengono introdotti i concetti fondamentali del platooning e le relative manovre, i principali controllori longitudinali (ACC/CACC) con attenzione ai requisiti di stabilità e ai meccanismi di degradazione, le tecnologie di comunicazione comunemente impiegate in ambito V2X (IEEE~802.11p, VLC, LTE/5G) e, infine, gli approcci multi-canale e le tecniche di selezione del link. L’analisi culmina nella discussione dei meccanismi di transizione sicura tra modalità di controllo e interfacce di comunicazione, tema centrale del lavoro di riferimento su cui si basa la parte sperimentale di questa tesi \cite{segata2023multi-technology}.
